\section{Introduction}

Every few years, a technology comes along that makes people declare everything before it obsolete.
When ChatGPT launched in late 2022, the reaction was swift: \emph{``Deep learning is dead. The future is all about large language models.''}\footnote{This sentiment was widespread in tech commentary following ChatGPT's release. It had been building for years: Gary Marcus's widely-read \emph{Nautilus} piece ``Deep Learning Is Hitting a Wall'' \cite{marcus2022wall} is a prominent pre-ChatGPT articulation of the view that deep learning had reached its limits, a claim sharply contested by Yann LeCun and others at the time.}

This paper argues that reaction got things exactly backwards.

Large language models \emph{are} deep learning.
They are not a replacement for it --- they are its most ambitious expression yet.
To understand why, we need to go back to the beginning.

Deep learning is, at its core, the science of teaching computers to learn from examples rather than explicit rules.
Instead of telling a program \emph{how} to recognize a cat, you show it ten million pictures of cats and let it figure out the patterns on its own.
The ``deep'' part refers to the many layers of processing through which information passes, each layer extracting more abstract patterns than the last.

This idea did not arrive overnight.
It took decades of failed experiments, stretches of neglect, and a handful of breakthrough moments that each changed the field's direction entirely.
Understanding that journey --- where deep learning came from, where it struggled, and what finally unlocked its potential --- makes the current moment in AI far more comprehensible.

So let's start at the beginning.
